% swim-times-conops.tex
% Concept of Operations for the swim-times program
% IEEE Std 1362-1998 structure
% Compiled with: pdflatex swim-times-conops.tex

\documentclass[12pt]{article}

\usepackage{geometry}
\geometry{letterpaper, margin=1in}

\usepackage{hyperref}
\usepackage{booktabs}
\usepackage{array}
\usepackage{enumitem}
\usepackage{titlesec}
\usepackage{fancyhdr}
\usepackage{xcolor}

% CAST blue
\definecolor{castblue}{RGB}{0,48,135}

\pagestyle{fancy}
\fancyhf{}
\renewcommand{\headrulewidth}{0.4pt}
\fancyhead[L]{\small College Area Swim Team}
\fancyhead[R]{\small ConOps: swim-times}
\fancyfoot[C]{\thepage}

\titleformat{\section}{\large\bfseries\color{castblue}}{}{0em}{}[\titlerule]
\titleformat{\subsection}{\normalsize\bfseries}{}{0em}{}

\newcommand{\code}[1]{\texttt{#1}}

\begin{document}

% -----------------------------------------------------------------
% Title block
% -----------------------------------------------------------------
\begin{titlepage}
  \centering
  \vspace*{2cm}
  {\Huge\bfseries\color{castblue} Concept of Operations\\[0.5em]}
  {\Large swim-times: Automated Swimmer-Times Retrieval System}\\[2em]
  {\large College Area Swim Team (CAST)}\\[0.5em]
  {\normalsize Prepared for: CAST Coaching Staff and Meet Directors}\\[0.5em]
  {\normalsize Prepared by: CAST Technology Working Group}\\[2em]
  \begin{tabular}{ll}
    Document identifier: & CAST-CONOPS-001                  \\
    Revision:            & 1.0                               \\
    Date:                & April 21, 2026                    \\
    Status:              & Released                          \\
    Governing standard:  & IEEE Std 1362-1998 (ConOps Guide) \\
  \end{tabular}
  \vfill
  {\small\textit{This document is prepared in accordance with the principles of
  rational documentation articulated in D.\,L.\,Parnas and P.\,C.\,Clements,
  ``A Rational Design Process: How and Why to Fake It,''
  \textit{IEEE Transactions on Software Engineering}, vol.\,SE-12, no.\,2,
  pp.\,251--257, February 1986.}}
\end{titlepage}

% -----------------------------------------------------------------
% Table of Contents
% -----------------------------------------------------------------
\tableofcontents
\clearpage

% =================================================================
\section{Scope}
% =================================================================

\subsection{System identification}

This document is the Concept of Operations (ConOps) for the
\textbf{swim-times} program, a compiled C utility that retrieves
competitive swim times from the USA~Swimming data hub and presents
them as formatted text or comma-separated-value (CSV) output.
The program is implemented as a CWEB literate program whose
canonical source is \code{swim2/stella-times.w}.

\subsection{Document purpose}

This ConOps establishes the operational rationale, usage model,
and stakeholder expectations for swim-times so that developers,
coaches, and meet directors share a common understanding of what
the system does and why it is built the way it is.
It is the primary reference for any subsequent requirements
specification or design document.

\subsection{System overview}

swim-times accepts command-line options that identify one or more
swimmers (Stella Julianna Evans, Kalea Rose Benavente, Kenneth Ray
Evans, and Keith Santiago Evans) and one or more USA~Swimming event
codes (e.g., \code{100 FR SCY}, \code{200 IM LCM}).  For each
selected swimmer--event pair it queries the live USA~Swimming
Sisense analytics back-end, retrieves all recorded times, sorts
them fastest-first, and prints either a human-readable table or a
CSV stream suitable for downstream processing.

% =================================================================
\section{Referenced Documents}
% =================================================================

\begin{enumerate}[label={[\arabic*]}]
  \item IEEE Std 1362-1998, \textit{IEEE Guide for Information Technology---
        System Definition---Concept of Operations (ConOps) Document}.
        IEEE Computer Society, 1998.

  \item D.\,L.\,Parnas and P.\,C.\,Clements, ``A Rational Design Process:
        How and Why to Fake It,''
        \textit{IEEE Transactions on Software Engineering},
        vol.\,SE-12, no.\,2, pp.\,251--257, February 1986.

  \item D.\,E.\,Knuth, ``Literate Programming,''
        \textit{The Computer Journal}, vol.\,27, no.\,2,
        pp.\,97--111, 1984.

  \item D.\,E.\,Knuth, \textit{The \TeX{}book}.
        Addison-Wesley, 1984.

  \item USA~Swimming, \textit{Individual Times Search},
        \texttt{https://data.usaswimming.org/datahub/usas/individualsearch}.
        Accessed 2026.

  \item Sisense Ltd., \textit{JAQL API Reference},
        \texttt{https://docs.sisense.com/main/SisenseLinux/jaql.htm}.
        Accessed 2026.

  \item \code{swim2/stella-times.w} --- CWEB literate program source
        (canonical source of record for the swim-times implementation).

  \item \code{swim/motivational-time-standards.pdf} --- USA~Swimming
        motivational time standards (B, BB, A, \ldots) used in
        the companion swim-times PDF report.

  \item \code{swim/sdi-agc-time-standards.pdf} --- San Diego-Imperial
        Age-Group Championship qualifying times.
\end{enumerate}

% =================================================================
\section{Current System / Situation (As-Is)}
% =================================================================

\subsection{Background}

CAST coaches and parents currently monitor swimmer progress by
visiting the USA~Swimming data hub at
\texttt{data.usaswimming.org}, entering each swimmer's name, and
manually reading and transcribing the resulting HTML tables.
Times are then typed into a Pandoc Markdown source file
(\code{swim/swim-times.md}) and compiled to PDF.

\subsection{Current problems}

\begin{description}[style=nextline]
  \item[Manual transcription errors.]
        Copying times by hand from a web browser introduces
        typographical mistakes; there is no automated verification
        that the entered time matches what USA~Swimming recorded.

  \item[Stale data.]
        The manual workflow is updated intermittently.  After any
        sanctioned meet new times may sit in the USA~Swimming
        database for days before the PDF is regenerated.

  \item[Scalability.]
        The team now tracks four swimmers across thirty-one events
        each, yielding up to 124 lookup--copy--paste operations per
        full refresh.

  \item[Swimcloud unreliability.]
        The alternative consumer site \texttt{swimcloud.com} is
        known to contain errors and is therefore excluded from the
        approved workflow.
\end{description}

\subsection{Operational environment}

The current workflow runs on a macOS workstation equipped with
\code{pandoc} and the \code{xelatex} PDF engine.  No persistent
server infrastructure exists; all work is ad hoc and human-driven.

% =================================================================
\section{Justification for Change}
% =================================================================

\subsection{Deficiencies addressed}

The swim-times program eliminates all four problems identified above:

\begin{enumerate}
  \item Live data is fetched directly from the authoritative
        USA~Swimming Sisense back-end, removing the human
        copy-paste step entirely.
  \item Times are always current because each invocation queries
        the live API; no stale cache is involved.
  \item Four swimmers and thirty-one events can be refreshed in a
        single command, reducing a multi-hour manual task to a few
        seconds of CPU time.
  \item Swimcloud is bypassed; the only source consulted is
        \texttt{data.usaswimming.org}.
\end{enumerate}

\subsection{Strategic alignment}

Automating time retrieval supports CAST's goal of making accurate,
up-to-date performance data available to coaches immediately after
every meet.  CSV output additionally enables integration with
spreadsheets, charting tools, and the team's existing Pandoc-based
PDF generation pipeline.

% =================================================================
\section{Proposed System (To-Be)}
% =================================================================

\subsection{System description}

swim-times is a single-binary command-line utility compiled from
the CWEB literate program \code{stella-times.w}.  It requires no
daemon, no database, and no network service beyond HTTPS access to
\texttt{usaswimming.sisense.com}.

\subsection{High-level capabilities}

\begin{description}[style=nextline]
  \item[Swimmer selection (\code{-o} flag).]
        The operator specifies one or more swimmer keywords
        (\code{stella}, \code{kalea}, \code{kenny}, \code{keith})
        to restrict output.  When no swimmer keyword is given all
        four swimmers are processed.

  \item[Event filtering (\code{-e} flag).]
        One or more USA~Swimming event codes restrict output to
        specific events.  The flag may be repeated or the codes
        may be comma-separated.  When \code{-e} is omitted all
        thirty-one events (fourteen SCY, seventeen LCM) are fetched.

  \item[Output modes.]
        \code{fastest} emits only the single fastest recorded time
        per event.  \code{csv} emits a CSV stream with a header
        row, suitable for import into spreadsheet software or the
        Pandoc PDF pipeline.

  \item[Live authoritative data.]
        All times are fetched from the USA~Swimming Sisense
        ElastiCube via authenticated JAQL API calls; no local cache
        is consulted.

  \item[Motivational standard annotation.]
        Each retrieved time includes the USA~Swimming motivational
        standard attained (B, BB, A, \ldots, or ``Slower Than B'')
        as recorded by USA~Swimming.
\end{description}

\subsection{Operational environment}

swim-times runs on any POSIX-compatible system with:

\begin{itemize}
  \item A C99 compiler (\code{gcc} or \code{clang}).
  \item \code{libcurl} with TLS support (for HTTPS to Sisense).
  \item CWEB (\code{ctangle}/\code{cweave}) if the literate source
        must be retangled or reweaved.
  \item \code{pdftex} if the typeset documentation PDF is desired.
\end{itemize}

\subsection{Interfaces}

\begin{description}[style=nextline]
  \item[Input.]
        Command-line arguments (\code{-o}, \code{-e}).

  \item[Network.]
        HTTPS POST to
        \texttt{https://usaswimming.sisense.com/api/datasources/}
        with a static bearer token embedded in the binary.

  \item[Standard output.]
        Formatted plain-text tables or CSV, depending on options.

  \item[Standard error.]
        Diagnostic messages only; no data is written to stderr.
\end{description}

% =================================================================
\section{Operational Scenarios}
% =================================================================

The following scenarios illustrate the primary use cases.

\subsection{Scenario 1 --- Post-meet fastest-times update}

\textit{Context:} It is the evening after a Saturday SCY meet.
A coach wants to check whether any of the four tracked swimmers
set new personal bests.

\begin{enumerate}
  \item Coach opens a terminal.
  \item Coach runs:
        \begin{center}
          \code{./swim-times -o stella,kalea,kenny,keith -o fastest}
        \end{center}
  \item swim-times resolves each swimmer's PersonKey from the
        USA~Swimming Persons table.
  \item For each swimmer, all thirty-one events are queried and
        the fastest single time is printed.
  \item Total elapsed time is typically under two minutes for all
        four swimmers.
  \item Coach reads the terminal output and identifies any new bests.
\end{enumerate}

\subsection{Scenario 2 --- CSV export for PDF report}

\textit{Context:} A parent wants to regenerate the landscape PDF
report (\code{swim/swim-times.pdf}) for Stella with current times.

\begin{enumerate}
  \item Parent runs:
        \begin{center}
          \code{./swim-times -o stella,csv > stella-latest.csv}
        \end{center}
  \item swim-times emits a CSV header followed by one row per
        recorded time for each of the thirty-one events.
  \item Parent (or a script) reads \code{stella-latest.csv}
        and updates \code{swim/swim-times.md}.
  \item Parent runs:
        \begin{center}
          \code{pandoc swim-times.md -o swim-times.pdf
                --pdf-engine=xelatex}
        \end{center}
  \item The updated PDF is ready.
\end{enumerate}

\subsection{Scenario 3 --- Single-event qualification check}

\textit{Context:} Kalea's coach wants to verify her current
100 Backstroke SCY time ahead of a championship entry deadline.

\begin{enumerate}
  \item Coach runs:
        \begin{center}
          \code{./swim-times -o kalea,fastest -e "100 BK SCY"}
        \end{center}
  \item swim-times resolves Kalea's PersonKey, queries the
        100~BK SCY event only, and prints the fastest time,
        date, meet, and motivational standard.
  \item Coach compares the printed time against the SDI~AGC
        qualifying standard (from
        \code{swim/sdi-agc-time-standards.pdf}) to determine
        eligibility.
\end{enumerate}

\subsection{Scenario 4 --- Full event history for coaching analysis}

\textit{Context:} Kenneth's coach wants to see his complete
200~Freestyle SCY progression to evaluate training effectiveness.

\begin{enumerate}
  \item Coach runs:
        \begin{center}
          \code{./swim-times -o kenny -e "200 FR SCY"}
        \end{center}
  \item swim-times retrieves all recorded 200~FR~SCY times for
        Kenneth, sorts them fastest-first, and prints a table
        showing time, date, standard, and meet name.
  \item Coach reviews the progression, noting improvement trends
        and peak performance periods.
\end{enumerate}

% =================================================================
\section{Summary of Impacts}
% =================================================================

\subsection{Staffing and workflow}

No dedicated operator is required.  Any team member with access
to the compiled binary and a network connection can invoke
swim-times.  The manual transcription step is eliminated,
reducing the time required for a full four-swimmer refresh from
approximately one to two hours to under five minutes.

\subsection{Training}

Users need only understand the two flags, \code{-o} and \code{-e},
and the swimmer and event keywords.  The built-in usage message
(printed when no arguments are supplied) provides sufficient
reference material for new users.  No formal training programme
is required.

\subsection{Infrastructure}

No new servers, databases, or persistent services are required.
The binary must be recompiled if the Sisense bearer token rotates
or if USA~Swimming changes its JAQL endpoint URLs.
These are maintenance events, not operational events.

\subsection{Security}

The Sisense bearer token is embedded at compile time.  Because
USA~Swimming's data hub serves only public competitive records,
the token grants read-only access to non-sensitive data.
Token rotation requires a source edit and recompile; no runtime
credential injection is provided.

\subsection{Compatibility}

The existing Pandoc-based PDF pipeline (\code{swim/swim-times.md})
is unchanged.  swim-times augments rather than replaces it,
supplying live data that a human or script pastes into the Markdown
source.  Future work may automate that final step as well.

% =================================================================
\section{Analysis of Alternatives}
% =================================================================

\subsection{Alternatives considered}

\begin{description}[style=nextline]
  \item[Web scraping via \code{curl}/HTML parsing.]
        Directly scraping the HTML rendered by
        \texttt{data.usaswimming.org} was considered.  It was
        rejected because the site renders content via JavaScript
        after an initial page load; a plain \code{curl} request
        returns an empty data table, making reliable extraction
        fragile and brittle.

  \item[Swimcloud API.]
        Swimcloud (\texttt{www.swimcloud.com}) aggregates
        USA~Swimming data but has known inaccuracies.  It was
        rejected in favour of the authoritative USA~Swimming source.

  \item[Full JSON parsing library (e.g., \code{cJSON}, \code{jansson}).]
        A third-party JSON parser would have provided a more robust
        response-processing layer.  It was rejected because the
        Sisense JAQL response format is stable and well-understood,
        and adding a library dependency would complicate the build
        for no operational gain.

  \item[Python or shell script.]
        A Python script using \code{requests} was prototyped.
        It was set aside in favour of a compiled C program to
        (a)~match the literate programming aesthetic of the
        CWEB toolchain already used in this project and
        (b)~avoid requiring a Python interpreter on the build host.
\end{description}

\subsection{Rationale for chosen approach}

The chosen design---a CWEB literate C program using libcurl and
string-scanning---satisfies all operational requirements with
minimal external dependencies and produces a self-documenting
artefact (the typeset \code{stella-times.pdf}) that serves
simultaneously as executable software and readable reference
documentation.  This directly embodies the Parnas--Clements
principle: the documentation and the implementation are the
same source file, presented in the order most useful to the
reader rather than the order dictated by the compiler.

\vfill
\begin{center}
  \small\textit{End of document CAST-CONOPS-001 Rev.\,1.0}
\end{center}

\end{document}
